% Notes of the Montreal seminar, 1970 (cote 129, typed part): archive cote 129, site page 8 (PDF page 9).
% typescript. Draft compared: none.
% AI-transcribed by Claude Opus 5.5 (claude-opus-5-5, adaptive thinking, effort high), Message Batches API, batch mtl129_full (msgbatch_01AGnDKHwy6ATCM7ehUa6263), 2026-11-13; not yet checked by eye.
\scanpage{129/8}
Comme $k$ est parfait, l'homomorphisme $u$ définit un isomor-%
phisme de groupes
\[ u(k) : \underset{\to}{\mathbb{W}}(k) \to \underset{\to}{\mathbb{W}}'(k) . \]

Remarquons alors que la projection naturelle $\pi_n : W \to W_n$
définit un isomorphisme
\[ \varepsilon_n : W/p^nW \xrightarrow{\ \sim\ } W_n \]
et l'on vérifie immédiatement que l'on a un diagramme commutatif
\diagram{
\begin{tikzcd}
W \arrow[rrr, "p\,=\,VF"] \arrow[dr] \arrow[dd, "\pi_n"'] & & & W \arrow[dl] \arrow[dd, "\pi_{n+1}"] \\
& W/p^nW \arrow[r, "p"] \arrow[dl, "\varepsilon_n", "\sim"'] & W/p^{n+1}W \arrow[dr, "\varepsilon_{n+1}"', "\sim"] & \\
W_n \arrow[rrr, "TF\,=\,FT"] & & & W_{n+1}
\end{tikzcd}
}

Utilisant alors l'isomorphisme naturel (de $W$-modules) valable pour tout
anneau de valuation discrète
\[ \underset{\to}{\lim}(W/p^nW,p) \xrightarrow{\ \sim\ } K/W \]
on déduit de ce qui précède que $\underset{\to}{\mathbb{W}}'(k)$ et donc $\underset{\to}{\mathbb{W}}(k)$ est isomorphe
au module dualisant $K/W$. Comme $K/W$ est un $W$-module, on déduit une
structure de $W$-module sur $\underset{\to}{\mathbb{W}}(k)$ par transport de structure.

\uline{Proposition} (4.2). Cette structure se prolonge de manière unique en
une structure de ind-schéma en groupes à anneau d'opérateurs $W$ sur
\[ \underset{\to}{\mathbb{W}}_k = \underset{\to}{\lim}((\mathbb{W}_n)_k , T) \]

\mspage{\hand{8}}
